\documentclass{article}
% \usepackage{enumitem}
\usepackage{enumerate}
\usepackage[english]{babel}
% \usepackage[utf8]{inputenc}
\usepackage{ctex}
% \usepackage[style=authoryear]{biblatex}
% \addbibresource{sample.bib}
% \newlist{steps}{enumerate}{1}
% \setlist[steps, 1]{label = Step \arabic*:}
\usepackage[letterpaper,top=2cm,bottom=2cm,left=2cm,right=2cm,marginparwidth=1.75cm]{geometry}

% Useful packages
\usepackage{nameref}
\usepackage{amsmath}
\usepackage{graphicx}
\usepackage{lipsum}
\usepackage{caption}
\usepackage{tikz}
\usepackage{subfig}
\usepackage{float}
\usepackage{overpic}
\usepackage[dvipsnames]{xcolor}
\usepackage{multirow}
%\usepackage{hyperref}
% \newcommand{\refSec}[1]{ \textcolor{blue}{\texttt{ \textbf{Yue:} #1~ }} }

% \newcommand{\refSec}[1]{ \textbf{Section:} \ref{#1}~ } 

\newcommand{\boldnameref}[1]{
    \textbf{\nameref{#1}}
}

\newcommand{\ResultFigWidth}{0.98\textwidth}
\newcommand{\gridMseFigWidth}{0.19\textwidth}

\newcommand{\GridMSE}[3]{results/GridScenes/Grid#1/Grid#2/alpha0.1/mse-#3.png}
% e.g. GridMSE{Pyramid}{25}{h2mc}

\newcommand{\GridFig}[3]{results/GridScenes/Grid#1/Grid#2/alpha0.1/result-#3.png}
\newcommand{\GridErrFig}[3]{results/GridScenes/Grid#1/Grid#2/alpha0.1/err-#3.png}
\newcommand{\PlaneErrFig}[3]{results/PlaneScenes/Plane#1/#2/err-#3.png}
% \newcommand{\VeachdoorFig}[3]{results/Veachdoor/alpha0.03-exponent100/result-h2mc.png}

\date{}
\title{实验报告}

\begin{document}

\maketitle
\indent 本报告首先文字性总结一下目前得到的结论(\boldnameref{sec:sum})。之后介绍一下实验场景与算法的默认配置(\boldnameref{sec:desc})。然后展示具体验证结论的rendering结果图片以及MSE曲线(\boldnameref{sec:benefit},\boldnameref{sec:cost})。之后会展示一些不同配置的情况,比如更改alpha,MLT算法,Chain数量等(\boldnameref{sec:otherConfig})。最后是一些我在实验过程中观察到的一些现象，以及一些想法，想和大家讨论一下(\boldnameref{sec:otherTopic})。
\section{回顾与总结}
\label{sec:sum}
\noindent \textbf{回顾：}\\
\indent 我们将MLT中路径贡献函数记作$f(x)$,其中$x$是primary space中的一点。通常在MLT中，我们一般将$f(x)$作为目标分布函数，实际上，我们也可以改变这一目标分布(新分布记作$g(x)$),同时将样本权重调整为($\frac{f(x)}{g(x)}$)，依然可以得到unbiased结果。对于包含sds路径的场景，即使是primary space下的$f(x)$，其有效区域依然是嵌入在全空间的低维流形。因此，我尝试将$g(x)$取为更加平坦的函数(如将$g(x)$取为场景中Specular物体的roughness放大后对应的路径贡献)。\textcolor{red}{$g(x)$的跳转接受率更高，但同时也会由于概率分布和真实贡献不同，引入额外的方差。（参见 \boldnameref{sec:cost}）}\\

\noindent \textbf{实验总结：}\\
\indent \textcolor{red}{在本阶段的实验中，我们发现随着场景中specular物体的增多，$g(x)$逐渐体现出更大的优势（MSE比$f(x)$的结果更低）。直观上的解释是：Specular物体的增多会导致完全不连通的sds pattern增多，遍历这些pattern只能通过large step跳转样本；而$g(x)$的large step接受率更高（Table \ref{tab:accept}），使得这些pattern之间的整体亮度比例可以更快地收敛。（实验验证 参见 \boldnameref{sec:benefit}）}

%另一方面，目标分布与实际的差异会引入额外的方差(可利用柯西积分不等式证明)。直观的解释是，在MH采样过程中部分样本stuck多次且$\frac{f(x)}{g(x)}$较大时，会放大误差，从而产生noise点。样本权重$\frac{f(x)}{g(x)}$较小时，也会降低稀释noise的能力。（参见 \boldnameref{sec:cost}）

\section{算法配置与实验场景}
\label{sec:desc}
\subsection{算法配置说明}
本次实验中透明材质为rough dielectric, 其NDF采用Beckman模型，粗糙度由$\alpha$描述；反射材质为phong模型(mala源码支持的模型)，粗糙度由exponent描述。本报告中默认$g_k(x)$表示将$\alpha$放大$k$倍，exponent缩小$k$倍后得到的路径贡献。由于MALA的源码中，direct path是直接渲染的，MLT只需要负责indirect path (bounce次数大于等于2)。因此，\textcolor{red}{本报告中所有图片都是render indirect path的图片，并将其normalize到相同平均亮度以排除不必要的干扰。}

\subsection{实验场景}
如无特殊标注，光源为一个很小的面光源，位于场景斜上方，本报告涉及到的场景如下:

\begin{enumerate}[(1)]
    \item \textbf{Grid[M]Pyramid}: Diffuse的地板上放置M个透明的金字塔，位置上排成网格，M=1,4,9,25,100。
    \item \textbf{Grid[M]Sphere}: Diffuse的地板上放置M个透明的圆球，位置上排成网格，M=1,4,9,25,100。
    \item \textbf{Grid[M]Teapot}: Diffuse的地板上放置M个透明茶壶(使用的是Veachdoor场景中的mesh)，位置上排成网格，M=1,4,9,25,100。
    % \item \textbf{Grid[M]TeapotRefl}: Diffuse的地板上放置M个phong材质的茶壶(使用的是Veachdoor场景中的mesh)，位置上排成网格，M=1,4,9,25,100。
    \item \textbf{PlaneSphere05}: Diffuse的地板上放置了半个表面有所起伏的透明球(05只是建模时试了不同起伏程度的编号，不用管)。
    \item \textbf{Slab}: SMS文中的场景。
    % \item \textbf{PlaneSlab}: 将Slab的曲面横置于Diffuse的地板上。
    \item \textbf{Veachdoor}: 同MALA源码中的场景，门缝外有一个较大的面光源。
    \item \textbf{Torus}: 同MALA源码中的场景，光源为环境光sunsky。
\end{enumerate}








\section{g(x)优点的验证}
\label{sec:benefit}
如Table \ref{tab:accept}所示，$g(x)$的接受率更高，尤其是large step较为显著，这使得不同sds区域之间的亮度比例可以更快收敛。Figure \ref{fig:Grid9PyramidFigMALA}、Figure \ref{fig:Grid100SphereFigH2MC}、Figure \ref{fig:Grid25TeapotFigMALA}分别以\textbf{Grid9Pyramid}、\textbf{Grid100Sphere}、\textbf{Grid25Teapot}为例展示了图片结果。通过error map可以更清晰地看出对于不同区域间亮度比例，$g(x)$收敛得更快，从而验证了上述分析。如Figure \ref{fig:GridPyramidMSE}、Figure \ref{fig:GridSphereMSE}、Figure \ref{fig:GridTeapotMSE}的MSE曲线所示，随着物体数量M的增加，场景中有更多完全不连通的sds区域，$g(x)$的优势更加明显。

\begin{table}[hb]
\centering
\begin{tabular}{|c|c|c|c|c|c|c|c|}
% \hline
% Grid25Pyramid-h2mc & l & l & l & s & s & s\\
\hline
Grid[M]Scene & M & f Large & $g_{1.25}$ Large & $g_{1.5}$ Large & f Small & $g_{1.25}$ Small & $g_{1.5}$ Small \\
\hline
\multirow{4}{*}{Grid[M]Pyramid} & 1 & 0.001351 & 0.002024 & 0.002766 & 0.5064 & 0.5132 & 0.5221\\
\cline{2-8} & 9 & 0.000988 & 0.00148 & 0.00204 & 0.3604 & 0.3631 & 0.3686\\
\cline{2-8} & 25 & 0.001205 & 0.001806 & 0.002492 & 0.290 & 0.2915 & 0.2958\\
\cline{2-8} & 100 & 0.001214 & 0.001809 & 0.002492 & 0.1608 & 0.1618 & 0.1642\\
\hline
\multirow{4}{*}{Grid[M]Sphere} & 1 & 0.0360 & 0.03873 & 0.04123 & 0.3946 & 0.4132 & 0.4339\\
\cline{2-8} & 9 & 0.02481 & 0.0267 & 0.02821 & 0.3589 & 0.3559 & 0.3574\\
\cline{2-8} & 25 & 0.02903 & 0.03135 & 0.03307 & 0.3324 & 0.3271 & 0.3255\\
\cline{2-8} & 100 & 0.02732 & 0.02964 & 0.03144 & 0.3016 & 0.2932 & 0.2878\\
\hline
\multirow{4}{*}{Grid[M]Teapot} & 1 & 0.00583 & 0.006402 & 0.006882 & 0.3478 & 0.3609 & 0.3704\\
\cline{2-8} & 9 & 0.006155 & 0.006661 & 0.007075 & 0.2439 & 0.2509 & 0.2547\\
\cline{2-8} & 25  & 0.007798 & 0.008409 & 0.008967 & 0.2169 & 0.2207 & 0.2224\\
\cline{2-8} & 100 & 0.007922 & 0.008557 & 0.009073 & 0.1880 & 0.1889 & 0.1879\\
\hline
% Grid9Pyramid-mala & 0.000987 & 0.00148 & 0.00204 & 0.2688 & 0.3298 & 0.3780\\
% \hline
% Grid100Sphere-mala & 0.02708 & 0.02941 & 0.03127 & 0.03293 & 0.04783 & 0.06356\\
% \hline
% Grid25Teapot-mala & 0.007857 & 0.008491 & 0.009032 & 0.06045 & 0.08629 & 0.1102\\
% \hline
\end{tabular}
\caption{Grid场景中，H2MC的接受率。其中$f$和$g_k$表示所使用的目标分布函数，其中$k$是$g(x)$对$\alpha$的放大倍率。Large和Small分别表示Large step和small step的平均接受率。以上数据为32768 spp测试下的平均结果。}
\label{tab:accept}
\end{table}

\clearpage
\begin{figure}
\begin{minipage}{\textwidth}
\centering  
\renewcommand{\arraystretch}{0.35}
\addtolength{\tabcolsep}{-5.0pt}
 \begin{tabular}{ ccccc }
\begin{overpic}[width=\gridMseFigWidth]{\GridMSE{Pyramid}{1}{h2mc}}\end{overpic}
& \begin{overpic}[width=\gridMseFigWidth]{\GridMSE{Pyramid}{4}{h2mc}}\end{overpic}
& \begin{overpic}[width=\gridMseFigWidth]{\GridMSE{Pyramid}{9}{h2mc}}\end{overpic}
& \begin{overpic}[width=\gridMseFigWidth]{\GridMSE{Pyramid}{25}{h2mc}}\end{overpic}
& \begin{overpic}[width=\gridMseFigWidth]{\GridMSE{Pyramid}{100}{h2mc}}\end{overpic}
\\
\begin{overpic}[width=\gridMseFigWidth]{\GridMSE{Pyramid}{1}{mala}}\end{overpic}
& \begin{overpic}[width=\gridMseFigWidth]{\GridMSE{Pyramid}{4}{mala}}\end{overpic}
& \begin{overpic}[width=\gridMseFigWidth]{\GridMSE{Pyramid}{9}{mala}}\end{overpic}
& \begin{overpic}[width=\gridMseFigWidth]{\GridMSE{Pyramid}{25}{mala}}\end{overpic}
& \begin{overpic}[width=\gridMseFigWidth]{\GridMSE{Pyramid}{100}{mala}}\end{overpic}
\\
\end{tabular}
\end{minipage}
\caption{Grid[M]PyramidMSE：Diffuse平面上摆放了M个透明金字塔，从左到右依次是M=1, 4, 9, 25, 100。第一行是h2mc的结果，第二行是mala的结果。图例中C128表示chain数量为128， P0.2表示large step的概率为0.2， roughScale表示$g(x)$对$\alpha$的放大程度，1.0表示$g(x)=f(x)$。从图中可以看到，随着M数量的增加，\textcolor{orange}{$g_{1.25}(x)$}和\textcolor{green}{$g_{1.5}(x)$}会优于\textcolor{blue}{$f(x)$}}
\label{fig:GridPyramidMSE} 
\end{figure}

\begin{figure}
\begin{minipage}{\textwidth}
\centering  
\addtolength{\tabcolsep}{-5.0pt}
 \begin{tabular}{ ccccc }
\begin{overpic}[width=\gridMseFigWidth]{\GridMSE{Sphere}{1}{h2mc}}\end{overpic}
& \begin{overpic}[width=\gridMseFigWidth]{\GridMSE{Sphere}{4}{h2mc}}\end{overpic}
& \begin{overpic}[width=\gridMseFigWidth]{\GridMSE{Sphere}{9}{h2mc}}\end{overpic}
& \begin{overpic}[width=\gridMseFigWidth]{\GridMSE{Sphere}{25}{h2mc}}\end{overpic}
& \begin{overpic}[width=\gridMseFigWidth]{\GridMSE{Sphere}{100}{h2mc}}\end{overpic}
\\
\end{tabular}
\end{minipage}
\caption{Grid[M]SphereMSE：Diffuse平面上摆放了M个透明玻璃球，从左到右依次是M=1, 4, 9, 25, 100。Small Step为H2MC。}
\label{fig:GridSphereMSE} 
\end{figure}

\begin{figure}
\begin{minipage}{\textwidth}
\centering  
\addtolength{\tabcolsep}{-5.0pt}
 \begin{tabular}{ ccccc }
\begin{overpic}[width=\gridMseFigWidth]{\GridMSE{Teapot}{1}{mala}}\end{overpic}
& \begin{overpic}[width=\gridMseFigWidth]{\GridMSE{Teapot}{4}{mala}}\end{overpic}
& \begin{overpic}[width=\gridMseFigWidth]{\GridMSE{Teapot}{9}{mala}}\end{overpic}
& \begin{overpic}[width=\gridMseFigWidth]{\GridMSE{Teapot}{25}{mala}}\end{overpic}
& \begin{overpic}[width=\gridMseFigWidth]{\GridMSE{Teapot}{100}{mala}}\end{overpic}
\\
% \begin{overpic}[width=\gridMseFigWidth]{\GridMSE{Teapot}{1}{h2mc}}\end{overpic}
% & \begin{overpic}[width=\gridMseFigWidth]{\GridMSE{Teapot}{4}{h2mc}}\end{overpic}
% & \begin{overpic}[width=\gridMseFigWidth]{\GridMSE{Teapot}{9}{h2mc}}\end{overpic}
% & \begin{overpic}[width=\gridMseFigWidth]{\GridMSE{Teapot}{25}{h2mc}}\end{overpic}
% & \begin{overpic}[width=\gridMseFigWidth]{\GridMSE{Teapot}{100}{h2mc}}\end{overpic}
% \\
\end{tabular}
\end{minipage}
\caption{Grid[M]TeapotMSE 在Diffuse平面上摆放了M个透明茶壶，从左到右依次是M=1, 4, 9, 25, 100。Small Step为MALA。}
\label{fig:GridTeapotMSE} 
\end{figure}
\begin{figure}
\begin{minipage}{\textwidth}
\centering  
\addtolength{\tabcolsep}{-5.0pt}
\begin{tabular}{c}
\begin{overpic}[width=\ResultFigWidth]{\GridFig{Pyramid}{9}{mala}}\end{overpic} \\
\begin{overpic}[width=\ResultFigWidth]{\GridErrFig{Pyramid}{9}{mala}}
    \put(14.3, 0){\begin{tikzpicture}[x=1pt,y=1pt]
        \draw[red,ultra thick] (0,0) rectangle (67.5,208);
    \end{tikzpicture}}
    \put(28.7, 0){\begin{tikzpicture}[x=1pt,y=1pt]
        \draw[red,ultra thick] (0,0) rectangle (67.5,208);
    \end{tikzpicture}}
    \put(43, 0){\begin{tikzpicture}[x=1pt,y=1pt]
        \draw[red,ultra thick] (0,0) rectangle (67.5,208);
    \end{tikzpicture}}
\end{overpic} \\
\end{tabular}
\end{minipage}
\caption{Grid9Pyramid-MALA: 第一行为$f(x)$的结果，第二行为$g_{1.25}(x)$，第三行为$g_{1.5}(x)$。后三行是对应的error map。通过error map的第二三四列都可以看出，$f(x)$的error map的区域间相对亮度误差更大。}
\label{fig:Grid9PyramidFigMALA}
\end{figure}

\begin{figure}
\begin{minipage}{\textwidth}
\centering  
\addtolength{\tabcolsep}{-5.0pt}
\begin{tabular}{c}
\begin{overpic}[width=\ResultFigWidth]{\GridFig{Sphere}{100}{h2mc}}\end{overpic} \\
\begin{overpic}[width=\ResultFigWidth]{\GridErrFig{Sphere}{100}{h2mc}}
    \put(7.4, 34.2){\begin{tikzpicture}[x=1pt,y=1pt]
        \draw[red, thick] (0,0) rectangle (10,10);
    \end{tikzpicture}}
    \put(21.8, 34.2){\begin{tikzpicture}[x=1pt,y=1pt]
        \draw[red, thick] (0,0) rectangle (10,10);
    \end{tikzpicture}}
    \put(36.2, 34.2){\begin{tikzpicture}[x=1pt,y=1pt]
        \draw[red, thick] (0,0) rectangle (10,10);
    \end{tikzpicture}}
    \put(38.7, 31.7){\begin{tikzpicture}[x=1pt,y=1pt]
        \draw[red, thick] (0,0) rectangle (10,10);
    \end{tikzpicture}}
    \put(50.6, 34.2){\begin{tikzpicture}[x=1pt,y=1pt]
        \draw[red, thick] (0,0) rectangle (10,10);
    \end{tikzpicture}}
    \put(53.1, 31.7){\begin{tikzpicture}[x=1pt,y=1pt]
        \draw[red, thick] (0,0) rectangle (10,10);
    \end{tikzpicture}}
    \put(65, 34.2){\begin{tikzpicture}[x=1pt,y=1pt]
        \draw[red, thick] (0,0) rectangle (10,10);
    \end{tikzpicture}}
    \put(11.3, 16.2){\begin{tikzpicture}[x=1pt,y=1pt]
        \draw[red, thick] (0,0) rectangle (10,10);
    \end{tikzpicture}}
    \put(25.7, 16.2){\begin{tikzpicture}[x=1pt,y=1pt]
        \draw[red, thick] (0,0) rectangle (10,10);
    \end{tikzpicture}}
\end{overpic} \\
\end{tabular}
\end{minipage}
\caption{Grid100Sphere-H2MC: 通过error map可以看出，$g_{1.25}(x)$和$g{1.5}(x)$相较于$f(x)$，各个区域亮度分配更均匀。}
\label{fig:Grid100SphereFigH2MC} 
\end{figure}

\begin{figure}
\begin{minipage}{\textwidth}
\centering  
\addtolength{\tabcolsep}{-5.0pt}
\begin{tabular}{c}
\begin{overpic}[width=\ResultFigWidth]{\GridFig{Teapot}{25}{mala}}\end{overpic} \\
\begin{overpic}[width=\ResultFigWidth]{\GridErrFig{Teapot}{25}{mala}}\end{overpic} \\
\end{tabular}
\end{minipage}
\caption{Grid25Teapot-MALA 这一场景中茶壶来自veachdoor中的mesh，从MSE来看，$g(x)$仍然优于$f(x)$的结果。}
\label{fig:Grid25TeapotFigMALA}
\end{figure}

\clearpage




\section{g(x)缺点的验证}
\label{sec:cost}
$g(x)$与真实贡献分布不一致会引入noise，并且随着$\alpha$的减小，$g(x)$带来的noise效应更加明显（如Figure \ref{fig:CostFigA0.1}($\alpha=0.1$)、Figure \ref{fig:CostFigA0.01}($\alpha=0.01$)、Figure \ref{fig:CostFigA0.005}($\alpha=0.005$)所示）。这里可以做一个近似的定量解释：

primaray space的$f(x) = \frac{L(path)}{P_{bdpt}(path)}$, 对于SDS路径, 实际的BDPT构造过程只能在其中一个Specular点上importance采样，另一个Specular点会直接连接Diffuse点（如ESD + SL或者ES+DSL）。无论哪一种，我们抛开一些平滑的项，近似得到$f(x)\approx BSDF(\theta)$，$\theta$为LightPath和EyePath连接处Specular点的法向与半角向量的夹角。进一步考虑NDF项，则有：

\begin{align}
\frac{f(x)}{g_k(x)} &\approx \frac{Beckman(\alpha, \theta)}{Beckman(k*\alpha, \theta)} \\
  &= \frac{(k*\alpha)^2}{\alpha^2} * e^{-\frac{tan^2\theta}{\alpha^2} + \frac{tan^2\theta}{(k*\alpha)^2}}\\
  &= k^2 * e^{-\frac{tan^2\theta}{\alpha^2} * (1 - \frac{1}{k^2})}
\end{align}

可以看出，对于相同的$g(x)$(即固定$k$)，$\alpha$越小，$\frac{f}{g}$随着$\theta$的变化越剧烈，导致概率分布和真实分布的差异越大，从而使得noise效应越明显。

我在Veachdoor和Torus中，也观察到了$g(x)$类似的缺陷（Figure \ref{fig:CostFigVeachdoor}, Figure \ref{fig:CostFigTorus}）。
\begin{figure}[ht]
\begin{minipage}[t]{0.1\textwidth}
\vspace{25pt}
\textbf{$f(x)$} \\
\\[33pt]
\textbf{$g_{1.25}(x)$} \\
\\[33pt]
\textbf{$g_{1.5}(x)$}
\end{minipage}%
\begin{minipage}[t]{0.9\textwidth}
\vspace{0pt}
\centering  
\begin{overpic}[width=\textwidth]{\PlaneErrFig{Pyramid}{alpha0.1}{h2mc}}\end{overpic}\\
\end{minipage}
\caption{Error Map of \textbf{PlanePyramid-H2MC, $\alpha=0.1$}}
\label{fig:CostFigA0.1} 
\end{figure}

\begin{figure}[ht]
\begin{minipage}[t]{0.1\textwidth}
\vspace{25pt}
\textbf{$f(x)$} \\
\\[33pt]
\textbf{$g_{1.25}(x)$} \\
\\[31pt]
\textbf{$g_{1.5}(x)$}
\\[40pt]
\textbf{$g_{2.0}(x)$}
\end{minipage}%
\begin{minipage}[t]{0.9\textwidth}
\vspace{0pt}
\centering  
\begin{overpic}[width=\textwidth]{\PlaneErrFig{Pyramid}{alpha0.01}{h2mc}}\end{overpic}\\
\end{minipage}
\caption{Error Map of \textbf{PlanePyramid-H2MC, $\alpha=0.01$}}
\label{fig:CostFigA0.01} 
\end{figure}

\begin{figure}[ht]
\begin{minipage}[t]{0.1\textwidth}
\vspace{25pt}
\textbf{$f(x)$} \\
\\[33pt]
\textbf{$g_{1.25}(x)$} \\
\\[31pt]
\textbf{$g_{1.5}(x)$}
\\[40pt]
\textbf{$g_{2.0}(x)$}
\end{minipage}%
\begin{minipage}[t]{0.9\textwidth}
\vspace{0pt}
\centering  
\begin{overpic}[width=\textwidth]{\PlaneErrFig{Pyramid}{alpha0.005}{h2mc}}\end{overpic}\\
\end{minipage}
\caption{Error Map of \textbf{PlanePyramid}, $\alpha=0.005$, H2MC.}
\label{fig:CostFigA0.005} 
\end{figure}

\begin{figure}[ht]
\begin{minipage}[t]{0.1\textwidth}
\vspace{10pt}
\textbf{$f(x)$} \\
\\[6pt]
\textbf{$g_{1.25}(x)$} \\
\\[5pt]
\textbf{$g_{1.5}(x)$}
\\[14pt]
\textbf{$g_{2.0}(x)$}
\end{minipage}%
\begin{minipage}[t]{0.9\textwidth}
\vspace{0pt}
\centering  
\begin{overpic}[width=\textwidth]{results/Veachdoor/alpha0.03-exponent100/result-h2mc.png}\end{overpic}\\
\end{minipage}
\caption{ \textbf{Veachdoor}, $\alpha=0.03, exponent=100$, H2MC.$g_{1.25}(x)$和$g_{1.5}(x)$的noise更大，尤其是最右侧的透明茶壶底部的pattern更加明显。相同spp下，MSE比$f(x)$更高。}
\label{fig:CostFigVeachdoor} 
\end{figure}

\begin{figure}[ht]
\begin{minipage}[t]{0.1\textwidth}
\vspace{15pt}
\textbf{$f(x)$} \\
\\[17pt]
\textbf{$g_{1.25}(x)$} \\
\\[14pt]
\textbf{$g_{1.5}(x)$}
\end{minipage}%
\begin{minipage}[t]{0.9\textwidth}
\vspace{0pt}
\centering  
\begin{overpic}[width=\textwidth]{results/Torus/alpha0.03-exponent100/result-h2mc.png}\end{overpic}\\
\end{minipage}
\caption{ \textbf{Torus}, $\alpha=0.03, exponent=100$, H2MC.相较于$f(x)$，$g_{1.25}(x)$和$g_{1.5}(x)$的noise更大。}
\label{fig:CostFigTorus} 
\end{figure}





\clearpage
\section{更多其他配置的实验}
\label{sec:otherConfig}
我们也测试了其他的一些情况，如不同Chain数量(Figure \ref{fig:ChainNumMSE})，不同的small mutation策略(Figure \ref{fig:MutationMSE})。将$g(x)$取为$g_k(x)$和$f(x)$的插值，并逐渐趋近于$f(x)$ (Figure \ref{fig:gxInterpolation})，这种情况下$g(x)$的MSE曲线会夹在$g_k(x)$和$f(x)$之间。这里如果是另一种思路，对large step使用$g_k(x)$，small step逐渐插值到$f(x)$，这么做会使得Markov的目标分布一直在$f$和$g$之间来回震荡，破坏了收敛性，数学上得不到正确的(unbiased)结果。

我们还测试了\textbf{Slab}，\textbf{PlaneSphere05}，\textbf{Grid25Sphere05}等，场景以及MSE如Figure \ref{fig:MoreScene}所示。
\begin{figure}
\begin{minipage}{\textwidth}
\centering  
\addtolength{\tabcolsep}{-5.0pt}
 \begin{tabular}{ ccc }
\begin{overpic}[width=0.3\textwidth]{results/PlaneScenes/PlanePyramid/alpha0.01/mse-chain_32-h2mc.png}\end{overpic}
& \begin{overpic}[width=0.3\textwidth]{results/PlaneScenes/PlanePyramid/alpha0.01/mse-chain_128-h2mc.png}\end{overpic}
& \begin{overpic}[width=0.3\textwidth]{results/PlaneScenes/PlanePyramid/alpha0.01/mse-chain_512-h2mc.png}\end{overpic}
\\
\end{tabular}
\end{minipage}
\caption{\textbf{PlanePyramid},$\alpha=0.01$，H2MC. 固定总样本数，改变chain数量，从左到右依次为32，128，512。}
\label{fig:ChainNumMSE}
\end{figure}

\begin{figure}
\begin{minipage}{\textwidth}
\centering  
\addtolength{\tabcolsep}{-5.0pt}
 \begin{tabular}{ ccc }
\begin{overpic}[width=0.3\textwidth]{results/PlaneScenes/PlanePyramid/alpha0.01/mse-pssMLT.png}\end{overpic}
& \begin{overpic}[width=0.3\textwidth]{results/PlaneScenes/PlanePyramid/alpha0.01/mse-chain_128-h2mc.png}\end{overpic}
& \begin{overpic}[width=0.3\textwidth]{results/PlaneScenes/PlanePyramid/alpha0.01/mse-mala.png}\end{overpic}
\\
\end{tabular}
\end{minipage}
\caption{\textbf{PlanePyramid},$\alpha=0.01$. 不同的small mutaion策略，依次为PSSMLT, H2MC, MALA}
\label{fig:MutationMSE}
\end{figure}


\begin{figure}
\begin{minipage}{\textwidth}
\centering  
\addtolength{\tabcolsep}{-5.0pt}
 \begin{tabular}{ c }
\begin{overpic}[width=0.7\textwidth]{results/PlaneScenes/PlanePyramid/alpha0.01/mse-roughinterpolation.png}\end{overpic}
\\
\end{tabular}
\end{minipage}
\caption{\textbf{PlanePyramid},$\alpha=0.01$. 图例中的后三种$g(x)$为$g_k(x)$到$f(x)$的插值。可以看出，$g(x)$的曲线在$g_k(x)$和$f(x)$(蓝色曲线)之间，k=1.25, 1.5, 2.0（这是早期的一个实验结果图，当时没有将图片亮度normalize，MSE的量级和本文其他地方不一致；且当时用了一个seed，因而曲线有些抖动）。}
\label{fig:gxInterpolation}
\end{figure}

\begin{figure}
\begin{minipage}{\textwidth}
\centering  
\renewcommand{\arraystretch}{0.35}
\addtolength{\tabcolsep}{-5.0pt}
\begin{tabular}{ cccc }
\begin{overpic}[width=0.24\textwidth]{results/Slab/alpha0.01/mse-h2mc.png}\end{overpic}
& \begin{overpic}[width=0.24\textwidth]{results/PlaneScenes/PlaneSphere05/alpha0.01/mse-h2mc.png}\end{overpic}
& \begin{overpic}[width=0.24\textwidth]{results/PlaneScenes/PlaneSphere05/alpha0.1/mse-h2mc.png}\end{overpic}
& \begin{overpic}[width=0.24\textwidth]{results/GridScenes/GridSphere05/Grid25/alpha0.1/mse-h2mc.png}\end{overpic}
\\
\begin{overpic}[width=0.24\textwidth]{results/Slab/alpha0.01/reference-pt.png}
\put(77, 60){\begin{tikzpicture}[x=1pt,y=1pt]
    \draw[red, thick] (0,0) rectangle (10,10);
\end{tikzpicture}}
\end{overpic}
& \begin{overpic}[width=0.24\textwidth]{results/PlaneScenes/PlaneSphere05/alpha0.01/reference-pt.png}\end{overpic}
& \begin{overpic}[width=0.24\textwidth]{results/PlaneScenes/PlaneSphere05/alpha0.1/reference-pt.png}\end{overpic}
& \begin{overpic}[width=0.24\textwidth]{results/GridScenes/GridSphere05/Grid25/alpha0.1/reference-pt.png}\end{overpic}
\\
\textbf{Slab}, $\alpha=0.01$ & \textbf{PlaneSphere05}, $\alpha=0.01$ & \textbf{PlaneSphere05}, $\alpha=0.1$ & \textbf{Grid25Sphere05}, $\alpha=0.1$
\end{tabular}
\end{minipage}
\caption{我们还测试了几个其他的场景，这里展示一下reference图以及对应的H2MC MSE曲线。其中Slab的MSE受到图中红色框内的一小块ES*L区域的影响，MSE很不稳定（也正因如此，在本报告中的场景设置中，我尽量避免了此类情况的发生）。\textbf{PlaneSphere05} $\alpha$越大, $\frac{f}{g}$越平坦，MSE曲线越接近；\textbf{Grid25Sphere05}的结果也说明，随着specular物体增多，SDS pattern增多，$g(x)$优势也会增大。}
\label{fig:MoreScene}
\end{figure}




\clearpage
\section{其他问题}

本小节是一些我在实验中的一些其他的失败尝试、观察到的现象以及一些想法。这里做一个简述，如果对更多具体结果感兴趣，可以会议上讨论或者我另外准备一份相关结果的汇报。

1. 我还尝试了一些其他的想法，对结果没有显著的提升，因此上文中没有提及。如(a)只对SDS路径使用$g(x)$。(b)proposal pdf仍然使用$f(x)$。(a)(b)的结果反而变差了。(c)Large step 目标分布使用$g(x)$, small step使用$f(x)$。(c)无法收敛到正确的结果（根据我自己的推导和资料查阅，理论上，两种不同目标分布的马尔科夫过程也难以线性组合，我没有得到组合后稳态分布的表达式）。

2. 在MALA的源码中，通过检测样本stuck达到一定次数后（我在本文实验中，使用的默认设置10000次），会强制将样本设为初始态，这样可以移除极端的noise。我尝试了将这个次数阈值减小，虽然可以更大程度上减少noise，但是这也会破坏Markov过程，使得结果有偏。

3.对于场景中包含ES*L路径的情况，往往会出现一小块过曝区域主导MSE，加上$\alpha$较小时，MSE很不稳定，不利于分析实验。

4.在实验中，我产生了一个与本项目无关的想法，这里提出来，向老师请教一下。对于Visibility很难搞的场景，比如光源通过小孔或者细缝照入场景，这使得哪怕是基于梯度的small step也很容易采到visibility无效的区域。然而实际上，在MLT过程中，有很多沿着细缝的成功样本可以利用。因此或许可以考虑利用这些样本矫正small step的proposal gaussian分布，从而使得mutation时以更高的概率沿着门缝采样。一种思路是，利用若干个历史样本估计当前样本primary space局部的visibility分布，将这一分布乘到原本的gaussian上。这种思路的潜在问题在于visibility未必容易刻画，刻画错了有时也可能带来额外cost。另一种思路是，想办法将visibility“软化”，得到visibility的近似微分，改进原有的梯度，得到包含visibility的梯度之后再proposal。

5.依然是一个与本项目无关的想法。假设Markov矩阵$M_1, M2, ...,M_n$都有相同的稳态分布$f$，易知对于$M=\sum_{i=1}^n \lambda_iM_i, where \sum_{i=0}^n \lambda_i = 1$，M的稳态分布依然是$f$。$\prod_{i=1}^n M_i$得到的新Markov矩阵也收敛到$f$。也就是说，可以将任意Proposal策略“并联”（线性组合）或者“串联”，仍然可以得到正确的MLT算法。LargeStep和SmallStep实际上就是两种Proposal的“并联”。同样的，将其他的MLT算法并联或者串联，或许可以平衡优劣，同时避免某种算法多次stuck形成noise点的几率。

\label{sec:otherTopic}
\end{document}
